% ===========================================================================
% Chapter 7: What Phase 1 establishes, and its limits
% ===========================================================================
\chapter{What Phase 1 Establishes}
\label{ch:establish}

This chapter states what the fourteen experiments establish and within which limits.
\emph{Measured} marks a claim resting on radiative comparisons, reproducible to 0.0017\%
of the change itself;
\emph{bounded} marks one resting on abundances, reproducible only to 2.5278\%
(Experiment~10).

\section{Established results}
\label{sec:est-results}

\textbf{The absorber term works, and its effect is located (measured).} The model's
absorber changes the radiation field in proportion to its amount, locally in wavelength,
and within one band of altitude, 69 to 93~km (Chapter~\ref{ch:exp17}).

\textbf{Placement is a first-order control (measured).} The same absorber has 7.34 times
more effect moved up than moved down, while changing its spectrum at a fixed placement is
worth a factor of 4.95, and spectral shape stops mattering once the absorber leaves the
responding layer (Experiments~7, 8 and~10).

\textbf{The chemistry, not the modeller, decides placement, and the closest gas-phase
candidate fails on it (measured).} At the profile its own chemistry gives it, ClS\sub{2}
leaves the response band and the radiation field becomes that of no absorber, to
0.001487\% (\texttt{id.l2\_ratio\_pct}; Experiment~9).

\textbf{No species the model can represent passes all three tests (calculated from
saved model output).} The
screen uses only saved model output and stored cross sections (Experiment~12).

\textbf{Two published candidates are reassessed (calculated from published data).} The
proposed iron and sulfur minerals cannot absorb enough at 375~nm even as pure solids,
and the published rejection of ferric chloride does not survive if ferric chloride absorbs
in the droplets as it does in hydrochloric acid (Experiments~13 and~14). Neither result uses the model, and neither identifies the
absorber.

\textbf{The chemical responses are small and mostly unresolved (bounded).} Apart from SO
and SO\sub{3} when the absorber is moved up, by 8.404\% and 4.718\% (keys
\texttt{res.e7a.SO.inband} and \texttt{res.e7a.SO3.inband}), the changes in the sulfur
reservoirs inside the band are at or below the chemistry floor, or above it by too little
to settle anything (SO in Experiments~3 and~4B), and are reported only as upper bounds.

\section{Limits}
\label{sec:limits}

\begin{enumerate}[leftmargin=*, itemsep=2pt]
  \item \textbf{One column, one sun angle.} The model is a single global average at a
    zenith angle of 58 degrees, with every photolysis rate halved for day and night
    (Table~\ref{tab:setup}). Latitude and local time are not represented, and every run
    is a steady state.
  \item \textbf{Low vertical mixing.} Every run uses the Bierson and Zhang mixing
    profile~\cite{bierson2020}, orders of magnitude below what has since been derived
    for the clouds from observations and models~\cite{dai2023kzz, lefevre2024}. The vertical wind is switched off. The
    placement verdict of Experiment~9 has therefore been tested at one mixing strength
    only; Chapter~\ref{ch:phase2} proposes to test it at others.
  \item \textbf{Simplified cloud droplets.} The cloud droplets enter the radiation only
    through fixed, fitted cross sections for two droplet modes; the full droplet
    calculation in the code stays switched off in every run
    (Section~\ref{sec:discovery}).
  \item \textbf{An empirical reference.} The empirical absorber's spectrum and profile
    were fitted to the observed brightness of Venus. Spectral rankings measure resemblance
    to that fit, not correctness (Experiments~8 and~11).
  \item \textbf{Only what the network contains.} The screen can judge only species the
    model tracks and holds cross sections for: 23 of 62. Ferric chloride and the iron
    minerals have no chemistry in the model and could not be placed by it.
  \item \textbf{A floor from one pair.} The chemistry floor comes from a single pair of
    replicate runs and is a lower bound on the true numerical scatter
    (Section~\ref{sec:conv-floor}).
\end{enumerate}

None of these limits weakens a comparison between two runs, which share them; they
limit how far the verdicts carry to the real atmosphere, and the first two are the
subject of Phase~2.
